\documentclass[11pt,a4paper]{article}
\usepackage[utf8]{inputenc}
\usepackage[T1]{fontenc}
\usepackage{amsmath,amssymb,amsthm}
\usepackage{geometry}
\geometry{margin=2.5cm}
\theoremstyle{plain}
\newtheorem{theorem}{Theorem}[section]
\newtheorem{proposition}[theorem]{Proposition}
\newtheorem{lemma}[theorem]{Lemma}
\theoremstyle{definition}
\newtheorem{definition}[theorem]{Definition}
\title{Soluciones Tests 41-45: Análisis}
\author{Mathematical AI Agent}
\date{\today}
\begin{document}
\maketitle
\section{Problemas}

% ============================================================
% TEST 41: Convergente implica acotada
% ============================================================

\begin{theorem}[Convergente implica acotada]
Sea $(a_n)$ una sucesión convergente en $\mathbb{R}$. Entonces $(a_n)$ es acotada.
\end{theorem}

\begin{proof}
Sea $\lim_{n\to\infty} a_n = a$ para algún $a \in \mathbb{R}$. Por definición de convergencia, para $\varepsilon = 1 > 0$ existe $N \in \mathbb{N}$ tal que para todo $n > N$:
\[
|a_n - a| < 1 \quad \Longrightarrow \quad |a_n| < |a| + 1.
\]
Definimos:
\[
M = \max\{|a_1|, |a_2|, \ldots, |a_N|, |a| + 1\}.
\]
Entonces, para todo $n \in \mathbb{N}$:
\begin{itemize}
    \item Si $n \leq N$, entonces $|a_n| \leq \max\{|a_1|, \ldots, |a_N|\} \leq M$.
    \item Si $n > N$, entonces $|a_n| < |a| + 1 \leq M$.
\end{itemize}
En ambos casos, $|a_n| \leq M$ para todo $n \in \mathbb{N}$. Por lo tanto, $(a_n)$ es acotada.
\end{proof}

% ============================================================
% TEST 42: Unicidad del límite
% ============================================================

\begin{theorem}[Unicidad del límite]
Si una sucesión $(a_n)$ converge, entonces su límite es único.
\end{theorem}

\begin{proof}
Supongamos que $\lim_{n\to\infty} a_n = a$ y $\lim_{n\to\infty} a_n = b$. Debemos demostrar que $a = b$.

Sea $\varepsilon > 0$ arbitrario. Por la definición de convergencia:
\begin{itemize}
    \item Existe $N_1 \in \mathbb{N}$ tal que para todo $n > N_1$: $|a_n - a| < \dfrac{\varepsilon}{2}$.
    \item Existe $N_2 \in \mathbb{N}$ tal que para todo $n > N_2$: $|a_n - b| < \dfrac{\varepsilon}{2}$.
\end{itemize}
Tomamos $N = \max\{N_1, N_2\}$. Para todo $n > N$, por la desigualdad triangular:
\[
|a - b| = |a - a_n + a_n - b| \leq |a_n - a| + |a_n - b| < \frac{\varepsilon}{2} + \frac{\varepsilon}{2} = \varepsilon.
\]
Dado que $\varepsilon > 0$ es arbitrario, se concluye que $|a - b| = 0$, es decir, $a = b$.
\end{proof}

% ============================================================
% TEST 43: Límite del producto
% ============================================================

\begin{theorem}[Límite del producto]
Si $(a_n)$ y $(b_n)$ son sucesiones convergentes en $\mathbb{R}$ con $\lim_{n\to\infty} a_n = a$ y $\lim_{n\to\infty} b_n = b$, entonces:
\[
\lim_{n\to\infty} a_n b_n = ab.
\]
\end{theorem}

\begin{proof}
Sean $\varepsilon > 0$ arbitrario. Como $(a_n)$ converge, es acotada por el Teorema 41, de modo que existe $K > 0$ tal que $|a_n| \leq K$ para todo $n \in \mathbb{N}$.

Consideramos la identidad:
\[
|a_n b_n - ab| = |a_n b_n - a_n b + a_n b - ab| = |a_n(b_n - b) + b(a_n - a)| \leq |a_n| \cdot |b_n - b| + |b| \cdot |a_n - a|.
\]
Por la convergencia de $(a_n)$ y $(b_n)$:
\begin{itemize}
    \item Existe $N_1$ tal que $|a_n - a| < \dfrac{\varepsilon}{2(|b| + 1)}$ para $n > N_1$.
    \item Existe $N_2$ tal que $|b_n - b| < \dfrac{\varepsilon}{2(K + 1)}$ para $n > N_2$.
\end{itemize}
(Se usa $|b|+1$ y $K+1$ en el denominador para evitar división por cero.)

Para $n > N = \max\{N_1, N_2\}$:
\[
|a_n b_n - ab| \leq K \cdot \frac{\varepsilon}{2(K + 1)} + |b| \cdot \frac{\varepsilon}{2(|b| + 1)} < \frac{\varepsilon}{2} + \frac{\varepsilon}{2} = \varepsilon.
\]
Por lo tanto, $\lim_{n\to\infty} a_n b_n = ab$.
\end{proof}

% ============================================================
% TEST 44: Límite del inverso
% ============================================================

\begin{theorem}[Límite del inverso]
Si $(a_n)$ es una sucesión convergente con $\lim_{n\to\infty} a_n = a$ y $a \neq 0$, entonces:
\[
\lim_{n\to\infty} \frac{1}{a_n} = \frac{1}{a}.
\]
\end{theorem}

\begin{proof}
Como $a \neq 0$, sea $\delta = \dfrac{|a|}{2} > 0$. Por la convergencia de $(a_n)$, existe $N_1 \in \mathbb{N}$ tal que para todo $n > N_1$:
\[
|a_n - a| < \delta = \frac{|a|}{2}.
\]
Por la desigualdad triangular inversa, para $n > N_1$:
\[
|a_n| = |a - (a - a_n)| \geq |a| - |a - a_n| > |a| - \frac{|a|}{2} = \frac{|a|}{2} > 0.
\]
En particular, $a_n \neq 0$ para todo $n > N_1$.

Sea $\varepsilon > 0$ arbitrario. Por la convergencia de $(a_n)$, existe $N_2 \in \mathbb{N}$ tal que para todo $n > N_2$:
\[
|a_n - a| < \frac{a^2 \varepsilon}{2}.
\]
Para $n > N = \max\{N_1, N_2\}$:
\[
\left|\frac{1}{a_n} - \frac{1}{a}\right| = \left|\frac{a - a_n}{a \cdot a_n}\right| = \frac{|a_n - a|}{|a| \cdot |a_n|} < \frac{|a_n - a|}{|a| \cdot \frac{|a|}{2}} = \frac{2|a_n - a|}{a^2} < \frac{2 \cdot \frac{a^2 \varepsilon}{2}}{a^2} = \varepsilon.
\]
Por lo tanto, $\lim_{n\to\infty} \dfrac{1}{a_n} = \dfrac{1}{a}$.
\end{proof}

% ============================================================
% TEST 45: Completeza de R (sucesiones de Cauchy)
% ============================================================

\begin{definition}[Sucesión de Cauchy]
Una sucesión $(a_n)$ en $\mathbb{R}$ es de Cauchy si para todo $\varepsilon > 0$ existe $N \in \mathbb{N}$ tal que para todo $m, n > N$:
\[
|a_m - a_n| < \varepsilon.
\]
\end{definition}

\begin{theorem}[Completeza de $\mathbb{R}$]
Toda sucesión de Cauchy en $\mathbb{R}$ es convergente.
\end{theorem}

\begin{proof}
Sea $(a_n)$ una sucesión de Cauchy en $\mathbb{R}$. Demostraremos que $(a_n)$ converge.

\textbf{Paso 1: $(a_n)$ es acotada.}
Tomando $\varepsilon = 1$, existe $N_0$ tal que $|a_m - a_n| < 1$ para todo $m, n > N_0$. En particular, para $n > N_0$:
\[
|a_n| \leq |a_n - a_{N_0+1}| + |a_{N_0+1}| < 1 + |a_{N_0+1}|.
\]
Definiendo $M = \max\{|a_1|, \ldots, |a_{N_0}|, 1 + |a_{N_0+1}|\}$, tenemos $|a_n| \leq M$ para todo $n$.

\textbf{Paso 2: Construcción de un punto límite.}
Por el Teorema de Bolzano-Weierstrass, toda sucesión acotada en $\mathbb{R}$ tiene una subsucesión convergente. Sea $(a_{n_k})$ una subsucesión convergente con:
\[
\lim_{k\to\infty} a_{n_k} = L.
\]

\textbf{Paso 3: La sucesión completa converge a $L$.}
Sea $\varepsilon > 0$. Como $(a_n)$ es de Cauchy, existe $N_1$ tal que:
\[
|a_m - a_n| < \frac{\varepsilon}{2} \quad \text{para todo } m, n > N_1.
\]
Como $a_{n_k} \to L$, existe $K$ tal que $n_K > N_1$ y:
\[
|a_{n_K} - L| < \frac{\varepsilon}{2}.
\]
Para todo $n > N_1$, por la desigualdad triangular:
\[
|a_n - L| = |a_n - a_{n_K} + a_{n_K} - L| \leq |a_n - a_{n_K}| + |a_{n_K} - L| < \frac{\varepsilon}{2} + \frac{\varepsilon}{2} = \varepsilon.
\]
Dado que $\varepsilon > 0$ es arbitrario, concluimos que $\lim_{n\to\infty} a_n = L$. Por lo tanto, toda sucesión de Cauchy en $\mathbb{R}$ converge.
\end{proof}

\end{document}
